%lezione 8 del 25 marzo 2026
%FINALMENTE LA PARTE BELLA
\newpage
\chapter{Formalismo dei \textit{Path Integral}}

Questa seconda parte del corso è dedicata al formalismo dei \textit{path integral} di Feynman:
si tratta di un framework equivalente a quello considerato finora, ma che ha i vantaggi di essere 
Lagrangiano, rendendo più semplice un'estensione relativistica, e di essere facilmente estendibile a 
una teoria con infiniti gradi di libertà, cioè a una teoria di campo.

Questa formulazione fu proposta da Feynman
\footnote{Articolo originale \hyperlink{https://web.archive.org/web/20200917091657/https://authors.library.caltech.edu/47756/1/FEYrmp48.pdf}{[link]}}
nel 1948, sulla base di alcune osservazioni di Dirac.

L'idea fondante è trovare la probabilità di transizione di una funzione d'onda $\psi$ tra due punti dello
spazio delle configurazioni e tra due tempi $\left(q_a,t_a\right)\rightarrow\left(q_b,t_b\right)$\footnote
{In questa parte degli appunti si denoteranno con $q,p$ l'insieme di tutte le coordinate canoniche del
sistema, sono da intendersi come $q^\mu,p^\mu$}:
in meccanica classica questo è semplice, in quanto esiste un'unica curva di moto $q(t)$ che collega i
due punti, mentre in meccanica quantistica, come si vedrà in seguito, è necessario considerare \textit{tutti}
i possibili cammini che collegano i due punti e integrare al variare di essi con diversi pesi.

\section{Il propagatore}

Il concetto fondamentale di questa formulazione è quello di \textbf{propagatore} (o kernel), che è
strettamente legato all'operatore evoluzione temporale $U(t_b,t_a)$.\\
Si consideri un sistema a $f$ gradi di libertà con hamiltoniana $H$ e uno stato $\ket{\psi(t_a)}$:
lo stato evoluto a un generico tempo $t_b>t_a$ è
$$\ket{\psi(t_b)}=U(t_b,t_a)\ket{\psi(t_a)}$$
con $U(t_b,t_a)$ che soddisfa
$$i\hbar\pdv[]{}{t_b}U = H_bU \leftrightarrow S_bU=0$$
dove $S_b:=i\hbar\pdv[]{}{t_b}-H_b$ è l'operatore differenziale di Schrödinger.

Scrivendo lo stato evoluto nello spazio delle posizioni, si ottiene
$$\bra{q_b}\ket{\psi(t_b)}:=\psi(q_b,t_b) = \bra{q_b}U(t_b,t_a)\ket{\psi(t_a)}$$
sfruttando la completezza delle posizioni $\ket{q}$ si ottiene
\begin{equation}
\psi(q_b,t_b) = \int d^fq\bra{q_b}U(t_b,t_a)\ket{q}\bra{q}\ket{\psi(t_a)} = 
\int d^fq\bra{q_b}U(t_b,t_a)\ket{q}\psi(q_a,t_a)
\label{eq:kernel_first}
\end{equation}
\textbf{Def.} Il propagatore di Schrödinger tra $(q_a,t_a)$ e $(q_b,t_b)$ è
$$K(q_b,t_b,q_a,t_a) = K(b,a) := \bra{q_b}U(t_b,t_a)\ket{q_a}$$
cioè l'elemento di matrice di $U(t_b,t_a)$ corrispondente alle posizioni $q_a$ e $q_b$.

\textbf{Oss.} Il valore di $\psi(q_b,t_b)$ dipende dal valore di $\psi(q,t_a)$ in \textit{tutto} lo spazio
fisico: il fronte d'onda a un tempo $t_b>t_a$ dipende da tutto il fronte d'onda precedente, convoluto con
una funzione che ne descrive l'evoluzione temporale; questo è molto simile al principio di Huygens-Fresnel
per la descrizione della propagazione delle onde nel tempo.
\newpage

\textbf{Oss.} Il propagatore $K(b,a)$ soddisfa l'equazione di Schrödinger, infatti
$$S_bU=0\leftrightarrow i\hbar\pdv[]{}{t_b}\bra{q_b}U\ket{q_a}-\bra{q_b}H_bU\ket{q_a}=0$$
inserendo la completezza delle posizioni nel secondo addendo si ottiene
$$\bra{q_b}H_bU\ket{q_a} = \int d^fq \bra{q_b}H\ket{q}\bra{q}U\ket{q_a} = 
\int d^fq H(q,t)\delta^f(q-q_b)\bra{q}U\ket{q_a} = H_bK(b,a)$$
e si trova che
$$i\hbar\pdv[]{}{t_b}K(b,a) - H_bK(b,a)=0 \leftrightarrow S_bK(b,a)=0$$
Si definisce inoltre \textit{propagatore ritardato} la funzione $G(b,a) = K(b,a)\Theta(t_b-t_a)$, necessaria
in quanto il propagatore non è ben definito per $t_b<t_a$.

\noindent\fbox{%
    \parbox{\textwidth}{%
    \textit{Ripasso:} funzioni di Green
    
    Dato un operatore differenziale $L_X$, la sua funzione di Green è una funzione $G(x)$ tale che 
    $$L_XG(x) = \delta(x)$$
    La proprietà fondamentale di $G(x)$ è
    $$L_Xf(x)=g(x)\leftrightarrow f(x) = (G*g)(x)$$
    dove $*$ è la convoluzione tra funzioni
    $$(G*g)(x) :=\int_{-\infty}^{\infty}dyG(y)g(x-y)$$
    }%
}

Si può dimostrare che $G(b,a)$ è la funzione di Green dell'operatore $S_b$:
\begin{align*}	
    S_bG&=S_bK(b,a)\Theta(t_b-t_a) = i\hbar\pdv[]{K}{t_b}\Theta + i\hbar K\pdv[]{\Theta}{t_b} -H_bK\Theta =\\
    &=H_bK\Theta+i\hbar K(b,a)\delta(t_b-t_a)-H_bK\Theta=i\hbar\bra{q_b}U(t_b,t_b)\ket{q_a}\delta(t_b-t_a)=\\
    &=i\hbar\delta^f(q_b-q_a)\delta(t_b-t_a)
\end{align*}

\textbf{Oss.} Il significato fisico del propagatore può essere visto considerando una particella confinata
in un punto $q$ al tempo $t=0$: la funzione d'onda evoluta nel punto $x$ al tempo $t>0$ sarà
$$\psi(x,t) = \int d^f\bar q K(x,t,\bar q,0)\delta^f(\bar q-q) = K(x,t,q,0) $$
\newpage
\subsection{Primo legame con l'azione}
Prima di ricavare esplicitamente il propagatore di una particella lbera, è comodo fare due piccoli conti.

Si supponga, senza perdere di generalità, che il propagatore associato a una hamiltoniana della forma
$H = T + V$ sia della forma
$$K(x,t,q,0) = e^{\frac{i}{\hbar}S(x,t,q,0)}$$
con $S$ funzione complessa: si impone l'equazione di Schrödinger per $K$
\begin{gather*}
i\hbar\pdv[]{K}{t}=-\pdv[]{S}{t}e^{\frac{i}{\hbar}S}\\
-\frac{\hbar^2}{2m}\pdv[2]{K}{x}+VK=e^{\frac{i}{\hbar}S}\left[-\frac{1}{\hbar^2}\left(\pdv[]{S}{x}\right)^2 +
\frac{i}{\hbar}\pdv[2]{S}{t}+V\right]
\end{gather*}
e si ottiene che
$$\frac{1}{2m}\left(\pdv[]{S}{x}\right)^2-\frac{i\hbar}{2m}\pdv[2]{S}{x}+V(x) = -\pdv[]{S}{t}$$
che, per $\hbar\rightarrow0$, è \textbf{l'equazione di Hamilton-Jacobi} per l'azione classica $S$.\\
Questa relazione inizia a far intuire che esiste un legame profondo tra l'azione $S$ e il propagatore.

Il secondo conto è ricavare il valore dell'azione classica $S$ sulla curva di moto della particella libera: l'hamiltoniana è $H=\frac{p^2}{2m}$, le equazioni del moto sono $m\ddot{q}=0$ e le curve di moto sono\\
$q(t) = vt+q_0$. Dati due istanti di tempo $t_2>t_1$, l'azione è
$$S_{cl}=\int_{t_1}^{t_2}dt\frac{1}{2}m\dot{q}^2=\frac{1}{2}mv^2(t_2-t_1)=\frac{m(q_2-q_1)^2}{2(t_2-t_1)}$$
dove $q_i = q(t_i)$.

\subsection{Il propagatore della particella libera 1D}

